\scanpage{131}
\uline{Proposition} Soit $Q\in Q^*$, et soit $\lambda\in C$ un nombre algébrique. Conditions équivalentes

(a) Pour tout conjugué $\lambda_\alpha$ de $\lambda$, on a $\lambda_\alpha\overline\lambda_\alpha=Q$ i.e. $|\lambda_\alpha|^2=\unsure{Q}$.

(b) \deleted{\unsure{Réciproquement}} Il existe un corps \inserted{de nbs} \uline{totalement réel} $K$, tel que \deleted{\unsure{on a}} $K(\lambda)$ soit égal à $K$ ou une extension \uline{quadratique} purement imaginaire de $K$, et on a $\lambda\overline\lambda=Q$ i.e. $|\lambda|^2=Q$.

(b') Comme (b), avec $K=Q(a)\deleted{\unsure{=Q(\lambda)}}$, où $a=\lambda+\overline\lambda=\lambda+\dfrac Q\lambda$.

(b'') Il existe un nombre algébrique \deleted{\unsure{a}} \uline{tot} \uline{réel} $a$, avec \deleted{\unsure{$a^2-4Q$ nul}} $4Q-a^2$ nul ou totalement positif, tel que $\lambda$ satisfasse une équation
\[
\lambda^2-a\lambda+Q=0
\]

(c) Soit $\mathbb K$ l'extension alg. tot. réelle maximale de $Q$, et $\mathbb L$ l'unique extension quadratique de $\mathbb K$, galoisienne sur $Q$. On a $\lambda\in\mathbb L$, et $\deleted{\unsure{K\lambda}}\lambda\overline\lambda=Q$ i.e. $|\lambda|^2=Q$.

\uline{Dém}

a)$\Longrightarrow$b'') Posons $a=\lambda+\overline\lambda=\lambda+\dfrac Q\lambda$. Je dis que $a$ \uline{est tot. réel}. En effet, les conjugués de $a$ sont les $\deleted{\unsure{\sigma}}a_\alpha=\deleted{\unsure{\sigma}}\lambda_\alpha+\dfrac Q{\deleted{\unsure{\sigma}}\lambda_\alpha}$, or par hypothèse $\lambda_\alpha\overline\lambda_\alpha=Q$ donc $a_\alpha=\lambda_\alpha+\overline\lambda_\alpha\in R$. D'autre part, $\lambda$ est solution de l'équation
\inserted{$(T-\lambda)(T-\overline\lambda)=0$ i.e.}
\[
\unsure{T^2|T}-aT+Q=0,
\]
et par suite $\lambda_\alpha$ est solution de l'équation $T^2-a_\alpha T+Q_{\unsure{\alpha}}=0$, \deleted{de plus} l'autre solution étant $\dfrac Q{\lambda_\alpha}$ \inserted{(qui peut être égale à $\lambda_\alpha$)}, \deleted{\unsure{donc}} comme $\lambda_\alpha\overline\lambda_\alpha=Q$, les deux solutions sont $\lambda_\alpha$ et $\overline\lambda_\alpha$, donc complexes conjuguées. Cela montre que $a_\alpha^2-4Q\le0$, \deleted{\unsure{donc}} \inserted{i.e.} $4Q-a_\alpha^2\ge0$, donc $4Q-a^2$ est tot positif.

b'')$\Longrightarrow$b') Il est évident que $\lambda$ est de degré au plus deux sur \inserted{le corps tot. réel} $Q(a)$, \deleted{\unsure{de plus comme}} \deleted{$\unsure{\lambda\in Q(a)}$} \inserted{et $Q(\lambda)\supset Q(a)$ [car $a=\lambda+\dfrac Q\lambda$].} \deleted{\unsure{car que}} Comme $a^2-4Q\le0$, les solutions de l'équation $T^2-aT+Q=0$ sont $\lambda$ et $\overline\lambda$ \inserted{(cf lemme plus bas)}, et on a $\lambda\overline\lambda=Q$, $\lambda+\overline\lambda=a$.
\mspage{130}
